\section{What Dogs Actually Signal: The Biological Substrate}
\label{sec:biology}

Before asking whether a machine can translate a dog, we must ask what there
is to translate. Canine communication is multimodal, graded, and heavily
context-dependent. It is also, crucially, \emph{asymmetric}: much of what a
dog ``says'' is addressed to humans, and the dog's receptive understanding of
human signals far exceeds its productive repertoire. This section reviews
the channels and the evidence that each carries decodable information.
Table~\ref{tab:channels} summarizes the landscape.

\subsection{Vocalizations: barks, growls, whines}
The domestic dog is unusually vocal for a canid, a trait widely interpreted
as a product of domestication and paedomorphosis. The ELTE Budapest
ethology group has built the most systematic body of evidence on what barks
encode. Pongr\'acz and colleagues showed that human listeners can classify
barks recorded in different behavioral situations above
chance~\cite{pongracz2005barks}, and that acoustic parameters---fundamental
frequency, tonality, and inter-bark interval---carry emotional information
that listeners reliably extract~\cite{pongracz2006acoustic}. Later work
established that barks also carry \emph{individual identity}: listeners can
discriminate dogs from bark acoustics alone~\cite{molnar2006discriminate},
and machine classifiers applied to thousands of barks from a kennel of
Hungarian sheepdogs recovered both context and individual at rates well
above chance, with ``stranger at the fence'' and ``fight'' contexts most
separable and ``play'' least. Growls are even more informative: Farag\'o et
al.\ demonstrated that human listeners judge a dog's size and aggressive
intent from growl formants, and that growls produced in food-guarding versus
play contexts are acoustically and perceptually
distinct~\cite{farago2017growls}. Whines and whimpers, associated with
separation and submission, complete a repertoire that is graded rather than
discrete.

The engineering literature has begun to automate exactly these judgments.
The ICASSP 2018 study ``What is my dog trying to tell me?'' framed bark
analysis as automatic recognition of \emph{context} and \emph{perceived
emotion}~\cite{icassp2018dogbark}, explicitly building on the ELTE
psychophysics. The important nuance---which we will return to in
Section~\ref{sec:eval}---is that the target label is the \emph{human-perceived}
emotion or the \emph{eliciting context}, not the dog's ``intended message.''
The model learns a mapping that humans already agree on; it does not recover
a hidden sentence.

\subsection{The face: DogFACS and the ``puppy dog eyes'' effect}
Canine facial behavior is codable with the same rigor as human facial
action coding. DogFACS, the Dog Facial Action Coding
System~\cite{waller2013dogfacs}, adapts Ekman and Friesen's FACS to the
musculature of the dog face, defining action units (e.g., AU101, the inner
brow raiser) and ear action descriptors. Using DogFACS, Kaminski et al.\
showed that dogs produce markedly more facial movement---especially the
inner brow raiser that enlarges the apparent eye---when a human is
\emph{attending} to them than when the human's back is turned, and that the
presence of food does not drive the effect~\cite{kaminski2017attention}.
This is a striking result: it suggests the ``puppy dog eyes'' expression is
a \emph{communicative} signal tuned to human attention, possibly shaped by
selection during domestication, rather than a mere reflex of internal
state. A 2024 comparative study further reported that domestication
\emph{constrains} the dog's facial expressiveness relative to
wolves~\cite{sci2024domestication}, complicating any simple ``dogs emote
freely'' narrative.

Automated facial analysis is now an active computer-vision subfield.
Explainable CNN pipelines recognize emotional states from canine facial
action patterns~\cite{sci2022explainable}, improved optimization schemes
push expression-recognition accuracy~\cite{sci2023cnnewoa}, and 2025 work
automates emotional-expression analysis end to end~\cite{sci2025automated}
and standardizes facial landmark detection~\cite{sci2025landmarks}. These
are genuine advances---but, again, their labels are human-coded affective
categories or DogFACS action units, not propositions.

\subsection{The body: posture, tail, and gaze}
A large fraction of everyday canine signaling is postural: ear position,
tail carriage and wag direction, weight distribution, piloerection, and
gaze. Classic work reported lateralized tail-wagging biases tied to
approach versus withdrawal motivation, and the broader ethogram of
dominance, submission, play (the play bow), and stress is well documented in
the clinical literature. Body signals are often \emph{more} diagnostic than
vocalizations, a point commercial vendors themselves concede: Traini's
founder has stated that audio alone is insufficient and that combining body
language and facial expression raises accuracy. This is an implicit
admission that the ``translation'' is really multimodal affect
classification. Olfactory signaling---urine marking, pheromonal
communication---is enormously important to dogs but essentially inaccessible
to current consumer sensors, and we flag it as a blind spot rather than a
solved channel.

\subsection{Reception: dogs understand more than they say}
Any honest account must separate what dogs \emph{produce} from what they
\emph{comprehend}. The receptive side is impressive. Rico, a border collie,
provided the first evidence of ``fast mapping''---inferring a novel word's
referent by exclusion---in \emph{Science}~\cite{kaminski2004rico}. Chaser,
another border collie, was documented retrieving 1{,}022 named objects and
responding to simple noun--verb combinations~\cite{pilley2011chaser}. These
are exceptional, intensively trained individuals, but they establish an
upper bound on referential comprehension that is far higher than the
productive repertoire. The asymmetry matters: a future ``dialogue'' system
may succeed by letting dogs \emph{receive} structured human input (as
soundboards do) long before it can \emph{read} structured canine output.

\subsection{Soundboards: the closest thing to canine ``speech''}
The most direct attempt to give dogs a productive symbolic channel is the
augmentative interspecies communication (AIC) soundboard: a grid of buttons
that each speak a recorded word when pressed. Popularized by speech
pathologist Christina Hunger and the Instagram dog Bunny, the phenomenon is
now under controlled study by Federico Rossano's UC San Diego
\emph{TheyCanTalk} project, which tracks thousands of button-using pets
across dozens of countries. Two 2024 peer-reviewed papers anchor the
science. Bastos et al.\ (\emph{PLoS ONE}) showed that soundboard-trained dogs
respond appropriately to words such as ``play'' and ``outside'' whether the
word is spoken by the owner or triggered by an unfamiliar button-presser,
indicating word comprehension rather than cue-reading~\cite{bastos2024plos}.
A companion \emph{Scientific Reports} paper analyzed two-button
combinations and found them to be non-accidental, non-random, and
non-imitative---evidence of intentional, if linear,
combination~\cite{bastos2024scirep}. A 2025 review synthesizes the evidence
and the skepticism~\cite{animalcog2025review}. The consensus interpretation
is careful: the combinations look like associative sequences (name\,+\,action,
place\,+\,go), not recursive grammar; spatial and immediate concepts
(``outside,'' ``paw'') are robust, while temporal abstractions
(``yesterday,'' ``tomorrow'') are rare and inconsistent. Soundboards thus
demonstrate \emph{referential intent} on the third rung of our ladder---but
they do so by giving the dog a human-designed symbol set, not by decoding
the dog's natural signal.

\begin{table}[htbp]
\centering
\caption{Canine communication channels, the information they carry, and the
strongest evidence. ``Decodable today'' refers to validated, above-chance
machine or human inference on a defined label---not to sentence-level
translation.}
\label{tab:channels}
\small
\begin{tabular}{@{}p{2.5cm}p{4.1cm}p{4.0cm}p{2.6cm}@{}}
\toprule
\textbf{Channel} & \textbf{Information encoded} & \textbf{Key evidence} & \textbf{Decodable today} \\
\midrule
Bark & Affect, arousal, context, individual identity &
Pongr\'acz et al.~\cite{pongracz2005barks,pongracz2006acoustic};
Moln\'ar et al.~\cite{molnar2006discriminate} & Affect/context (above chance) \\
Growl & Body size, aggressive intent, resource-guarding vs.\ play &
Farag\'o et al.~\cite{farago2017growls} & Size/intent (human + ML) \\
Whine/whimper & Separation distress, submission & Ethological observation &
Coarse affect \\
Face (DogFACS) & Attention-directed signals, affective state &
Waller et al.~\cite{waller2013dogfacs}; Kaminski et al.~\cite{kaminski2017attention} &
Action units; affect (CNN) \\
Body/tail/gaze & Motivation, stress, play, social intent &
Clinical ethograms; wag-laterality studies & Posture classes \\
Olfaction & Identity, reproductive state, territory &
Scent-marking literature & \emph{Not} sensor-accessible \\
Reception (input) & Human word referents (up to 1{,}022) &
Kaminski et al.~\cite{kaminski2004rico}; Pilley \& Reid~\cite{pilley2011chaser} &
Comprehension, not production \\
Soundboard (AIC) & Referential intent via human symbols &
Bastos et al.~\cite{bastos2024plos,bastos2024scirep} &
Word-level intent (trained) \\
\bottomrule
\end{tabular}
\end{table}

\subsection{The takeaway}
The biological substrate is real and rich: dogs signal affect, arousal,
identity, motivation, and---with human scaffolding---referential intent.
What the substrate does \emph{not} contain is an open-ended generative
grammar. This single fact, more than any limitation of current models,
bounds what ``translation'' can mean. We formalize that bound next.